% Generated draft: requires mathematical and source review.
\documentclass[11pt]{article}
\usepackage[margin=25mm]{geometry}
\usepackage{fontspec}
\setmainfont{Latin Modern Roman}
\usepackage{amsmath,amssymb,mathtools,unicode-math}
\setmathfont{Latin Modern Math}
\usepackage{enumitem,needspace,xurl,hyperref}
\setlength{\parindent}{0pt}
\setlength{\parskip}{5pt}
\setlength{\emergencystretch}{3em}
\newcommand{\N}{\mathbb{N}}
\newcommand{\Z}{\mathbb{Z}}
\newcommand{\Q}{\mathbb{Q}}
\newcommand{\R}{\mathbb{R}}
\newcommand{\C}{\mathbb{C}}
\newcommand{\F}{\mathbb{F}}
\newcommand{\A}{\mathbb{A}}
\newcommand{\PP}{\mathbb{P}}
\newcommand{\E}{\mathbb{E}}
\newcommand{\eps}{\varepsilon}
\newcommand{\dd}{\,\mathrm{d}}
\begin{document}
\section*{UnsolvedMath ID: 800613}
\textbf{Source problem code:} catalogue entry
\subsection{Definitions and mathematical statement}
Let \(G\) be a separable metrizable topological group with identity \(e\). Here separability means that \(G\) has a countable dense subset, and metrizability refers to its topology. The supplied excerpt does not define its cover notation. The following gives a standard interpretation; its agreement with the source requires review, particularly concerning the exclusion of covers containing \(G\) itself.

Let \(\mathcal O(G)\) be the family of open covers \(\mathcal U\) of \(G\) such that \(G\notin\mathcal U\). Let \(\Omega(G)\subseteq\mathcal O(G)\) consist of the \(\omega\)-covers: those covers for which every finite subset of \(G\) is contained in some member of the cover. For an open neighborhood \(U\) of \(e\), put
\[
\mathcal O(U)=\{gU:g\in G\},\qquad
\Omega(U)=\{FU:\varnothing\ne F\subseteq G\text{ is finite}\},
\]
where \(gU=\{gu:u\in U\}\) and \(FU=\bigcup_{g\in F}gU\). Define
\[
\begin{aligned}
\mathcal O_{\mathrm{nbd}}(G)
 &=\{\mathcal O(U):U\text{ is an open neighborhood of }e,\ U\ne G\},\\
\Omega_{\mathrm{nbd}}(G)
 &=\{\Omega(U):U\text{ is an open neighborhood of }e,
       \ FU\ne G\text{ for every finite }F\subseteq G\}.
\end{aligned}
\]
Thus the members of \(\Omega_{\mathrm{nbd}}(G)\) are \(\omega\)-covers under the stated convention. Left translates are used here; the excerpt does not specify the translation convention.

For families \(\mathcal A\) and \(\mathcal B\) of covers of \(G\), the property \(\mathsf S_1(\mathcal A,\mathcal B)\) means that for every sequence \((\mathcal U_n)_{n\ge1}\) from \(\mathcal A\), there exist choices \(V_n\in\mathcal U_n\) such that \(\{V_n:n\ge1\}\in\mathcal B\). The property \(\mathsf S_{\mathrm{fin}}(\mathcal A,\mathcal B)\) means that for every such sequence there exist finite subfamilies \(\mathcal V_n\subseteq\mathcal U_n\) such that \(\bigcup_{n\ge1}\mathcal V_n\in\mathcal B\). Empty finite subfamilies are permitted. If \(\mathcal A\) is empty, these universally quantified properties hold vacuously.

Write
\[
\begin{aligned}
P(G)&:\quad\mathsf S_{\mathrm{fin}}\bigl(\mathcal O_{\mathrm{nbd}}(G),\mathcal O(G)\bigr),\\
Q(G)&:\quad\mathsf S_1\bigl(\Omega_{\mathrm{nbd}}(G),\Omega(G)\bigr).
\end{aligned}
\]
The source asks whether these properties determine the same class of separable metrizable groups:
\[
\forall G\;\bigl[G\text{ is a separable metrizable topological group}
\ \Longrightarrow\ (P(G)\Longleftrightarrow Q(G))\bigr].
\]
It specifically asks the following two set-theoretic questions.
\begin{enumerate}
\item Does the Continuum Hypothesis \(2^{\aleph_0}=\aleph_1\) imply that there exists a separable metrizable topological group \(G\) such that \(P(G)\) holds and \(Q(G)\) fails?
\item Is it consistent that every separable metrizable topological group \(G\) satisfies \(P(G)\Longleftrightarrow Q(G)\)?
\end{enumerate}
The usual intended background for these questions is \(\mathsf{ZFC}\), but the supplied excerpt does not explicitly specify a base theory or a relative-consistency hypothesis. Equality here means equality of classes of groups satisfying the properties. The supplied registry labels the problem open; this draft makes no independent assertion about its present status.
\subsection{Short English statement}
The problem compares two selection properties of separable metrizable topological groups: making finitely many choices from each neighborhood-translate cover to cover the group, and making one choice from each finite-translate neighborhood cover so that every finite subset is contained in a chosen set. It asks whether the Continuum Hypothesis yields a group with the first property but not the second, and whether equality of the two classes is consistent. The excerpt omits the underlying cover conventions, so the explicit interpretation above requires source review.
\subsection{Sources}
\begin{itemize}
\item Local UnsolvedMath registry, permanent ID \(800613\), exact source problem code \texttt{catalogue entry}; supplied complete source record, AMR Open Problem Lists. The record gives upstream revision \texttt{d371ee917ac9647fbb06cd6fc529deec6f3dda82} and labels the problem \emph{open}. This is supplied registry metadata, not an independently verified mathematical status.
\item Tsaban, \emph{Selection Principles and Special Sets of Reals: Open Problems} (2006), Problem 41, \url{https://arxiv.org/abs/math/0609601}. The registry lists Tsaban as proposer and 2006 as the proposed year. Its background records extraction by \texttt{arxiv\_latex} and an access date of 2026-10-16; these are supplied provenance statements, not evidence of access or verification during this draft.
\item The supplied attribution/context is the reference label \texttt{[BG]}. Its bibliographic expansion and precise attribution role are absent from the supplied record and remain unresolved.
\end{itemize}
\end{document}
